% ---
% title: Monte Callo Method
% date: 2026/1/2
% tag: ML
% ---
\section{Monte Callo Method}

Monte Callo Method is a method for analyzing the uncertainty of a model.

\section{Introduction}

设\textbf{随机变量\(X,Y\)}：\(Y=f(X),\quad X\sim p(x)\)

\textbf{真实期望}：\(\mathbb E[Y]=\mathbb E[f(X)]=\int_{-\infty}^{\infty} f(x)p_X(x)dx\)

\textbf{蒙特卡洛估计}：\(\boxed{ \hat{\mu}=\frac{1}{N}\sum_{i=1}^{N} f(x_i), \quad x_i\sim p(x) }\)

\section{蒙特卡洛不确定性分析}

蒙特卡洛做“不确定性分析”的核心就是：\textbf{把不确定的输入当作随机变量抽样很多次→每次跑一遍模型→统计输出分布}，从而得到输出的不确定区间、超标概率、敏感性等。

\subsection{1）先把问题写成：输出 = 模型(输入)}

假设你关心一个指标 \(Y\)（比如检测精度、成本、风险、产量…），它由一堆输入 \(X_1,\dots,X_d\) 决定：

\[
Y = f(X_1,\dots,X_d)
\]

不确定性就体现在 \(X_i\) 不是固定值，而是“可能取值的分布”。

\hrule

\subsection{2）给每个不确定输入建分布（最关键的一步）}

常见做法：

\begin{itemize}
  \item \textbf{有历史数据}：直接拟合分布（或用经验分布/核密度）
  \item \textbf{只有上下界}：常用 \textbf{均匀分布 Uniform(a,b)} 或 \textbf{三角分布 Triangular(min, mode, max)}
  \item \textbf{只有均值±方差}：常用 \textbf{正态 Normal(μ,σ)}（注意物理下界时用截断正态）
  \item \textbf{比例/概率型变量（0\textasciitilde{}1）}：常用 \textbf{Beta}
  \item \textbf{正数且偏态}（时间、乘性误差）：常用 \textbf{LogNormal}
  \item \textbf{离散选择}（策略A/B/C）：用分类分布（按概率抽）
\end{itemize}

如果输入之间有关联（相关系数/共变结构），要同时建\textbf{相关性}（不然会低估或高估不确定性）。

\hrule

\subsection{3）抽样 N 次：随机生成 N 组输入}

最基础：\textbf{独立随机抽样（i.i.d.）}。

更推荐（收敛更快）：

\begin{itemize}
  \item \textbf{拉丁超立方抽样 LHS}
  \item \textbf{Sobol / Halton 低差异序列}（准蒙特卡洛）
\end{itemize}

\hrule

\subsection{4）逐次运行模型，得到 N 个输出}

对 \(k=1..N\)：

\[
x^{(k)} \sim p(X),\quad y^{(k)} = f(x^{(k)})
\]

最终得到输出样本 \(\{y^{(k)}\}\)，这就是“输出的不确定性”。

\hrule

\subsection{5）输出你想要的不确定性结论}

常见指标（按你论文/报告里最常用的）：

\begin{itemize}
  \item \textbf{分布}：直方图/核密度
  \item \textbf{均值、方差、分位数}：比如 P5/P50/P95
  \item \textbf{置信区间/可信区间}（在仿真语境里通常直接报“分位区间”更直观）
  \item \textbf{超标/失效概率}：\(\Pr(Y>阈值)\)
  \item \textbf{贡献分析/敏感性}（见下一节）
\end{itemize}

\hrule

\subsection{6）做“谁导致不确定性最大”的分析（敏感性）}

两类最常用：

\subsubsection{A. 便宜快：相关/回归类}

\begin{itemize}
  \item \textbf{PRCC（部分秩相关）}：非线性但单调时很好用
  \item \textbf{标准化线性回归系数（SRC）/秩回归（SRRC）}
\end{itemize}

输出：每个 \(X_i\) 对 \(Y\) 的影响排序。

\subsubsection{B. 更严谨：Sobol 全局敏感性}

把输出方差分解：

\begin{itemize}
  \item \textbf{一阶 Sobol} \(S_i\)：单独变量贡献
  \item \textbf{总效应} \(S_{T_i}\)：包含交互项的总贡献
\end{itemize}

输出：能看出“交互”是否重要。

\hrule

\subsection{7）\(N\) 取多大？（实用规则）}

\begin{itemize}
  \item 如果只是看均值/分位数粗略范围：\textbf{\(N=10^3 \thinspace 10^4\)}
  \item 需要稳定估计尾部概率（很小的失效概率）：可能要 \textbf{更大 N} 或用重要性抽样等
  \item 你可以做一个简单的\textbf{收敛检查}：N 翻倍，看 P95/失效概率是否变化很小（比如 <1\% 相对变化）
\end{itemize}

\hrule

\subsection{8）最常见坑}

\begin{itemize}
  \item \textbf{分布乱选}：输出看着很“科学”，其实输入假设错了就全错
  \item \textbf{忽略相关性}：会显著改变尾部风险
  \item \textbf{把“参数不确定”与“随机性/噪声”混在一起}：建议分层（见下）
  \item \textbf{模型运行太慢}：用代理模型（高斯过程/随机森林/多项式混沌）或并行
\end{itemize}

\subsection{9）进阶：分层蒙特卡洛（区分两种不确定性）}

\begin{itemize}
  \item \textbf{Aleatory（随机性）}：不可约（环境波动）
  \item \textbf{Epistemic（认知不确定）}：可通过更多数据减少（参数未知）
\end{itemize}

常用做法：外层抽参数（epistemic），内层抽噪声（aleatory），最后给出“区间里的区间”。

\hrule

如果你给我一句话描述你的场景（比如：你模型输出是什么、哪些输入不确定、有没有数据/上下界），我可以直接按你的变量类型给一套\textbf{分布选择 + 抽样方案（LHS/Sobol）+ 输出统计 + 敏感性}的模板（偏论文写法那种）。
